\section*{Hoofdstuk \ref{ch:g12:diffeq} --- Differentiaalvergelijkingen}

\begin{solution}{exo:g12:diffeq:1}
\emph{(a)} $y = C\eu^{3x}$; $y(0) = 2$ geeft $y = 2\eu^{3x}$.

\emph{(b)} $y' = -\frac12 y$, dus $y = C\eu^{-x/2}$; $y(0) = -1$ geeft
$y = -\eu^{-x/2}$.

\emph{(c)} Evenwicht $y_p = 5$; $y = C\eu^{-x} + 5$; $y(0) = 0$ geeft
$C = -5$: $y = 5\left(1 - \eu^{-x}\right)$.
\end{solution}

\begin{solution}{exo:g12:diffeq:2}
$N' = kN$, dus $N(t) = N_0 \eu^{kt}$. Verdubbelen in $3$ uur betekent
$\eu^{3k} = 2$, dus
\[
k = \frac{\ln 2}{3} \approx 0.231\ \text{u}^{-1}.
\]
\end{solution}

\begin{solution}{exo:g12:diffeq:3}
$y_p'(x) = \eu^x + x\eu^x = y_p(x) + \eu^x$: $y_p$ is een particuliere
oplossing (het resonante geval van \cref{met:g12:diffeq:particular}). Volgens
\cref{prop:g12:diffeq:general} zijn de oplossingen
$y = C\eu^{x} + x\eu^{x} = (C + x)\,\eu^x$ met $C \in \R$.
\end{solution}

\begin{solution}{exo:g12:diffeq:4}
Probeer $y_p = \beta\eu^x$: uit $\beta\eu^x = -2\beta\eu^x + \eu^x$ volgt
$3\beta = 1$, dus $y_p = \frac13\eu^x$. Algemene oplossing
$y = C\eu^{-2x} + \frac13\eu^{x}$; de voorwaarde $y(0) = 1$ geeft
$C = \frac23$:
\[
y(x) = \frac{2}{3}\,\eu^{-2x} + \frac{1}{3}\,\eu^{x}.
\]
\end{solution}

\begin{solution}{exo:g12:diffeq:5}
$N(t) = N_0\,\eu^{-\lambda t}$ met $\lambda = \frac{\ln 2}{5730}$
(\cref{ex:g12:exp:halflife}). We lossen $\eu^{-\lambda t} = 0.6$ op:
\[
t = \frac{\ln(1/0.6)}{\lambda} = 5730\,\frac{\ln(5/3)}{\ln 2}
\approx 5730 \times \frac{0.5108}{0.6931} \approx 4220 \text{ jaar}.
\]
\end{solution}

\begin{solution}{exo:g12:diffeq:6}
\emph{1.} Er stroomt $0.2 \times 5 = 1$ kg zout per minuut binnen. De
uitstroom voert een concentratie $\frac{m}{100}$ kg/L af met $5$ L/min, dus
$\frac{m}{20}$ kg/min. Bijgevolg is $m' = 1 - \frac{m}{20}$.

\emph{2.} Evenwicht $m_p = 20$; $m(t) = 20 + C\eu^{-t/20}$, en $m(0) = 0$
geeft $C = -20$:
\[
m(t) = 20\left(1 - \eu^{-t/20}\right) \xrightarrow[t\to+\infty]{} 20 \text{ kg}.
\]
Op lange termijn is de concentratie in de tank gelijk aan die van de
binnenstromende pekel: $0.2$ kg/L $\times$ $100$ L $= 20$ kg.
\end{solution}

\begin{solution}{exo:g12:diffeq:7}
\emph{1.} $v' = g - \frac{k}{m} v$ met $\frac km = \frac{16}{80} = 0.2$.
Evenwicht $v_\infty = \frac{mg}{k} = \frac{9.8}{0.2} = 49$ m/s;
$v(t) = 49 + C\eu^{-0.2t}$, en $v(0) = 0$ geeft
\[
v(t) = 49\left(1 - \eu^{-0.2 t}\right).
\]

\emph{2.} Eindsnelheid $49$ m/s ($\approx 176$ km/u). We willen
$1 - \eu^{-0.2t} = 0.95$, dus $\eu^{-0.2t} = 0.05$:
\[
t = \frac{\ln 20}{0.2} \approx \frac{3.00}{0.2} \approx 15 \text{ s}.
\]
\end{solution}

\begin{solution}{exo:g12:diffeq:8}
\emph{1.} $z = \frac1y$ geeft
$z' = -\frac{y'}{y^2} = -\frac{y(1-y)}{y^2} = -\frac{1-y}{y}
= -\frac1y + 1 = -z + 1$.

\emph{2.} Evenwicht $z_p = 1$, dus $z(t) = 1 + C\eu^{-t}$. Uit
$y(0) = \frac{1}{10}$ volgt $z(0) = 10$, dus $C = 9$ en
\[
y(t) = \frac{1}{1 + 9\,\eu^{-t}} .
\]

\emph{3.} Voor $t \to +\infty$ is $9\eu^{-t} \to 0$ en $y(t) \to 1$. De
kromme is de klassieke S-vormige (\emph{sigmoïde}) logistische kromme: eerst
trage groei ($y$ klein, $y' \approx y$), de snelste groei bij $y = \frac12$
(waar $y' = y(1-y)$ maximaal is), en daarna verzadiging naar de
draagkracht $1$.
\end{solution}

\begin{solution}{exo:g12:diffeq:9}
Volgens \cref{thm:g12:diffeq:affine} is $y(x) = C\eu^{ax} - \frac ba$.
Bijgevolg is
\[
u_{n+1} = C\eu^{a(n+1)} - \frac ba
= \eu^{a}\left(C\eu^{an} - \frac ba\right) + \frac ba\left(\eu^a - 1\right)
= q\,u_n + r
\]
met $q = \eu^{a}$ en $r = \frac{b}{a}\left(\eu^{a} - 1\right)$. De voorwaarde
$\abs q < 1$ betekent $\eu^a < 1$, dus $a < 0$: precies de voorwaarde
waaronder de oplossingen van de \omterm{def:g12:diffeq:ode}{differentiaalvergelijking} naar het evenwicht
$-\frac ba$ convergeren --- en inderdaad is het vaste punt van de recursie
$\frac{r}{1 - q} = -\frac{b}{a}$.
\end{solution}

\begin{solution}{pb:g12:diffeq:1}
\textbf{1.} $y = 2\eu^{3t}$. Voor de tweede: evenwicht $3$, dus
$y = 3 + C\eu^{-2t}$ met $y(0) = 0$, waaruit $C = -3$:
$y = 3\left(1 - \eu^{-2t}\right)$.

\textbf{2.} $\left(y\eu^{-at}\right)' = y'\eu^{-at} - ay\eu^{-at} =
(y' - ay)\eu^{-at} = 0$: het product is een constante $C$, dus
$y = C\eu^{at}$ --- geen enkele oplossing ontsnapt.

\textbf{3.} Evenwicht: $y \equiv 3$. Elke andere oplossing is
$3 + C\eu^{-2t}$ met $C \neq 0$: de exponentiële term sterft uit en de
oplossing glijdt naar $3$ --- een stabiel \omterm{pb:g10:functions:1}{vast punt}, de continue tweelingbroer
van de vaste punten van de recursies uit \cref{pb:g11:seq:1}.

\textbf{4.} $y = y_0\eu^{-t/\tau}$ halveert wanneer $\eu^{-t/\tau} =
\frac12$: $t = \tau\ln 2$.

\textbf{5.} Alle oplossingen zijn verticale verschuivingen van het verval
naar $3$: die erboven starten zakken naar $3$, die eronder starten stijgen
naar $3$, en de constante oplossing $3$ blijft liggen --- een trechter rond
het evenwicht.

\textbf{6.} $\lambda = \frac{\ln 2}{5730} \approx 1.21 \times 10^{-4}$ per
jaar.

\textbf{7.} $\eu^{-\lambda t} = 0.2$: $t = \frac{\ln 5}{\lambda} \approx
13\,300$ jaar.

\textbf{8.} $t = \frac{\ln(1/0.15)}{\lambda} \approx 15\,700$ jaar: de
stieren van Lascaux dateren uit de late ijstijd.

\textbf{9.} $t = \frac{\ln(1/0.53)}{\lambda} \approx 5\,250$ jaar: binnen
een mensenleven van de waarde van de archeologen --- de klok werkt.

\textbf{10.} Tien halveringstijden laten $2^{-10} \approx 0.1\,\%$ over: aan
de rand van het meetbare. Na $65$ miljoen jaar is de fractie
$2^{-65\,000\,000/5\,730} \approx 2^{-11\,344}$: van de oorspronkelijke
voorraad blijft in geen enkel fossiel één atoom over --- dinosauriërs worden
met tragere klokken gedateerd (kalium-argon en verwanten).

\textbf{11.} $52\,\%$ geeft ongeveer $5\,405$ jaar en $54\,\%$ ongeveer
$5\,095$: de $\pm 1\,\%$ van de scheikunde wordt ruwweg $\pm 155$ jaar
geschiedenis --- nauwkeurigheid in het labo is nauwkeurigheid op het bordje
in het museum.

\textbf{12.} $z = T - T_a$ voldoet aan $z' = -kz$:
$z = (T_0 - T_a)\eu^{-kt}$, en $T = T_a + z$.

\textbf{13.} Het beginverschil is $T_0 - T_a = 70$; na vijf minuten is het
\omterm{def:g11:seq:arithmetic}{verschil} $70 - 20 = 50$, dus $50 = 70\eu^{-5k}$:
$k = \frac{\ln(70/50)}{5} \approx 0.067$ per minuut. Drinkbaar:
$55 - 20 = 35 = 70\eu^{-kt}$, dus $t = \frac{\ln 2}{k} \approx 10.3$
minuten.

\textbf{14.} De snelheid van het verlies is evenredig met het \omterm{def:g11:seq:arithmetic}{verschil}
$T - T_a$: gloeiend hete koffie bloedt het snelst warmte weg. De melk
\emph{nu} toevoegen snijdt het \omterm{def:g11:seq:arithmetic}{verschil} meteen weg, zodat er over de tien
minuten minder warmte verloren gaat; ze op het einde toevoegen laat de koffie
eerst op volle snelheid afkoelen. Voor de warmste kop: eerst de melk.
(Ongeduldige gastheren hebben de thermodynamica omgekeerd.)

\textbf{15.} Verschillen met de omgeving: om middernacht $10$, een uur later
$8$: $\eu^{-k} = 0.8$, dus $k = \ln\frac{10}{8} \approx 0.223$ per uur. Bij
het overlijden was het \omterm{def:g11:seq:arithmetic}{verschil} $17$; het zakte tot $10$ tegen middernacht:
$s = \frac{\ln(17/10)}{k} \approx 2.4$ uur. Tijdstip van overlijden: rond
21.40 uur.

\textbf{16.} Een verwarmde of tochtige kamer ($T_a$ niet constant) buigt de
klok naar beide kanten; kleding of lichaamsmassa verandert $k$ (de tabellen
van de wetsdokter corrigeren daarvoor) --- een verkeerde $k$ schaalt het hele
\omterm{def:g10:numbers:interval}{interval}; en een lichaam dat elders vandaan komt zet $T_a$ halverwege het
verval opnieuw, wat een vroeger of later overlijden voorspiegelt. De
\omterm{def:g10:algebra:equation}{vergelijking} is eerlijk; het risico zit in de aannames.

\textbf{17.} $z = \frac1y$ voldoet aan $z' = 1 - z$: $z = 1 + 9\eu^{-t}$
(uit $z(0) = 10$), dus $y = \frac{1}{1 + 9\eu^{-t}}$. \omterm{def:g12:deriv:inflection}{Buigpunt} in
$y = \frac12$: $9\eu^{-t} = 1$, dus $t = \ln 9 \approx 2.2$ --- de keerdatum
van de epidemie, rechtstreeks uit de formule.

\textbf{18.} Eindsnelheid: $v' = 0$ in $v = 20$ m/s. Oplossing:
$v(t) = 20\left(1 - \eu^{-t/2}\right)$. Dan $0.95$: $\eu^{-t/2} = 0.05$, dus
$t = 2\ln 20 \approx 6$ s --- regendruppels halen hun kruissnelheid binnen
enkele seconden, en daarom is regen niet dodelijk.

\textbf{19.} Evenwicht: $c = 8$. De helft ervan:
$4 = 8\left(1 - \eu^{-t/2}\right)$, dus $\eu^{-t/2} = \frac12$ en
$t = 2\ln 2 \approx 1.4$ uur. Het infuus is de continue tweelingbroer van de
lening: een constante instroom en een evenredige uitstroom ---
\cref{exo:g12:diffeq:9} maakt het woordenboek exact.

\textbf{20.} Verval ($a < 0$, $b = 0$), afkoeling ($y = T - T_a$), vallen
($a < 0$, $b > 0$: het evenwicht van onderaf naderen) en doseren (hetzelfde,
in een ader): één \omterm{def:g10:algebra:equation}{vergelijking}, vier werelden. De lus: formuleer de wet van
verandering; schrijf de \omterm{def:g10:algebra:equation}{vergelijking} op; los ze op
(\cref{met:g12:diffeq:solve}); ijk de constanten op metingen; voorspel; en
controleer daarna de aannames (vraag 16). Een trilling vraagt
$y'' = -\omega^2 y$ --- ontmoet, opgelost en aan het schommelen gebracht in
\cref{pb:g12:trigo:1}.
\end{solution}
